\subsection{室}

定义 2. 相对于 H 的 E 的一个室（如果关于 H 没有歧义，也简称为室）是相对于 H 的 E 的一个不包含于属于 H 的任何超平面的面片。

设 U 是 E 中由不属于 H 的任何超平面的点组成的开集；由于 H 中的超平面若与某个面片相交就必须包含该面片，室就是包含于 U 的面片；由命题 3 (i)，每个室都是 E 的凸（因而连通）开子集；由于室构成 U 的一个划分，它们恰好是 U 的连通分支。U 的每个凸子集 A 都是连通的，因而包含于一个室中；若 A 非空，该室是唯一的。显然，室就是支撑为 E 的面片，而命题 3 (iii) 表明每个室都是其闭包的内部。最后，设 C 是一个室，A 是 C 的非空子集；公式 (2) 和 (3) 蕴含

\[
C = \bigcap_ {H \in \mathfrak {H}} D _ {H} (A) = D _ {\mathfrak {H}} (A), \quad \overline {{C}} = \bigcap_ {H \in \mathfrak {H}} \overline {{D _ {H} (A)}}\tag{6}
\]

因为对于所有 \(D_{\mathrm{H}}(\mathbf{A}) = D_{\mathrm{H}}(a)\)，有 \(a \in A\)。

命题 5. 设 C 是 E 的非空子集。假设 \(\mathfrak{H}'\) 存在一个子集 \(\mathfrak{H}\)，具有下列性质：

\begin{enumerate}
\def\labelenumi{\alph{enumi})}
\item
  对于任意 \(\mathrm{H} \in \mathfrak{H}'\)，存在由 \(\mathrm{D}_{\mathrm{H}}\) 界定的开半空间 \(\mathrm{H}\)，使得 \(\mathrm{C} = \bigcap_{\mathrm{H} \in \mathfrak{H}'} \mathrm{D}_{\mathrm{H}}\)。
\item
  集合 C 不与属于 \(\mathfrak{H} - \mathfrak{H}'\) 的任何超平面相交。
\end{enumerate}

在这些条件下，C 是 \(\mathfrak{H}\) 在 E 中定义的一个室，并且对于所有 \(\mathrm{D_H} = \mathrm{D_H(C)}\)，\(\mathrm{H} \in \mathfrak{H}\)。

性质 a) 和 b) 表明 C 是 U 的凸子集；因此存在一个室 \(C'\)，使 \(C \subset C'\)。由于 \(C \subset D_{H}\)，对 \(D_{H} = D_{H}(C)\) 中的所有 H，有 \(H'\)，从而 \(C = D_{H'}(C) \supset D_{H}(C)\)，因为 \(H' \subset H\)；由 (6)，有 \(D_{H}(C) = C'\)，因而 \(C \supset C'\)。最终 \(C = C'\)。

命题 6. E 的每一点都至少属于一个室的闭包。

如果 E 退化为单点，这是显然的。否则，设 \(a \in E\)，令 \(H_{1}, \ldots, H_{m}\) 为 H 中包含 a 的超平面。由于 H 局部有限，存在 a 的一个邻域 V，它不与 H 中除 \(H_{1}, \ldots, H_{m}\) 之外的任何超平面相交。令 D 为经过 a 且不包含于任何 \(H_{i}\) 的直线；若 \(x \in D\)、\(x \neq a\)，且 x 充分接近 a，则开线段 \(|ax|\) 包含于 V 中，且不与任何 \(H_{i}\) 相交。于是 \(|ax| \subset U\)；由于 \(|ax|\) 连通，它包含于一个室 C 中，因而 \(a \in \overline{C}\)。

命题 7. 设 L 是 E 的一个仿射子空间，\(\Omega\) 是 L 的一个非空开子集。

\begin{enumerate}
\def\labelenumi{(\roman{enumi})}
\item
  \(a\) 中存在一点 \(\Omega\)，不属于 \(\mathfrak{H}\) 中不包含 \(L\) 的任何超平面。
\item
  如果 \(\mathrm{L}\) 是超平面且 \(\mathrm{L} \notin \mathfrak{H}\)，则存在一个与 \(\Omega\) 相交的室。
\item
  如果 \(\mathrm{L}\) 是超平面且 \(\mathrm{L} \in \mathfrak{H}\)，则 \(a\) 中存在一点 \(\Omega\)，不属于 \(\mathrm{H} \neq \mathrm{L}\) 中任何超平面 \(\mathfrak{H}\)。
\end{enumerate}

记 N 为满足 \(H \in H\) 且 \(L \notin H\) 的超平面 H 所成的集合，记 L 为仿射空间 L 中形如 \(L \cap H\)（其中 \(H \in N\)）的超平面所成的集合。显然，L 是 L 中的局部有限超平面集，而命题 6 表明 \(\Omega\) 与 L 在 L 中定义的一个室 \(\Gamma\) 相交。若 a 是 \(\Gamma \cap \Omega\) 中的一点，则对所有 \(a \notin H\) 有 \(H \in N\)，从而得到 (i)。

现在假设 L 是超平面；于是包含 L 的任意超平面都等于 L，因此可分为两种情形：

\begin{enumerate}
\def\labelenumi{\alph{enumi})}
\item
  \(\mathbf{L} \notin \mathfrak{H}\)：则 \(\mathfrak{N} = \mathfrak{H}\)，并且对于所有 \(a \notin \mathbf{H}\) 有 \(\mathbf{H} \in \mathfrak{H}\)；因此 \(a\) 属于 \(\mathfrak{H}\) 在 \(\mathbf{E}\) 中定义的一个室，从而得到 (ii)。
\item
  \(\mathrm{L} \in \mathfrak{H}\)：则 \(\mathfrak{N} = \mathfrak{H} - \{\mathrm{L}\}\)，从而得到 (iii)。
\end{enumerate}

